\begin{table}[H]
\centering
\caption{Standardized Effect Sizes for Main Outcomes}
\label{tab:sde}
\begin{threeparttable}
\begin{adjustbox}{max width=\textwidth}
\begin{tabular}{lcccccc}
\toprule
Outcome & $\hat{\beta}$ & SE & SD($Y$) & SDE & SE(SDE) & Classification \\
\midrule
\multicolumn{7}{l}{\textit{Panel A: Pooled}} \\
ln(Enterprises) --- Preferred (unit trends) & 0.080 & 0.029 & 2.17 & 0.037 & 0.013 & Small positive \\
ln(Enterprises) --- Na\"ive TWFE & -0.190 & 0.034 & 2.17 & -0.088 & 0.016 & Moderate negative \\
ln(Enterprises) --- Short window (2013+) & -0.091 & 0.024 & 2.17 & -0.042 & 0.011 & Small negative \\
ln(Enterprises) --- CZ only & -0.040 & 0.063 & 2.06 & -0.019 & 0.031 & Small negative \\
\midrule
\multicolumn{7}{l}{\textit{Panel B: Heterogeneous}} \\
ln(Enterprises) --- Cash-intensive sectors & -0.226 & 0.039 & 2.17 & -0.104 & 0.018 & Moderate negative \\
ln(Enterprises) --- Non-cash-intensive sectors & -0.184 & 0.038 & 2.17 & -0.085 & 0.017 & Moderate negative \\
\bottomrule
\end{tabular}
\end{adjustbox}
\par\vspace{0.3em}
{\footnotesize
\begin{itemize}[leftmargin=*]
\item \textit{Notes:} \textbf{Country:} Czech Republic. \textbf{Research question:} Does mandatory real-time electronic reporting of cash transactions affect the number of registered business enterprises in targeted economic sectors? \textbf{Policy mechanism:} The Electronic Records of Sales (EET) system required all businesses with cash transactions to transmit receipt data to the tax authority in real time, creating a technological enforcement tool that increased the cost of operating informally while potentially reducing the viability of marginal enterprises dependent on tax evasion. \textbf{Outcome definition:} Number of active enterprises at the 2-digit NACE division level, from Eurostat Structural Business Statistics, measured annually. \textbf{Treatment:} Binary indicator equal to one for Czech NACE divisions in EET-covered sectors after their respective phase-in date (NACE I and G from 2017; H, M, and C from 2018). \textbf{Data:} Eurostat SBS, 2008--2020, 2-digit NACE division $\times$ country $\times$ year panel with five Visegrad+ countries (CZ, SK, PL, HU, AT), balanced panel of 324 units. \textbf{Method:} Two-way fixed effects with division$\times$country and year fixed effects; preferred specification adds unit-specific linear trends. Standard errors clustered at division$\times$country level. Sun--Abraham event study as alternative heterogeneity-robust estimator. \textbf{Sample:} 2-digit NACE divisions with complete data 2008--2020; sectors B through S excluding agriculture (A, not in SBS) and public administration (O, P, Q excluded as non-market). SDE $= \hat{\beta} / \text{SD}(Y)$ where SD($Y$) is the unconditional standard deviation of ln(enterprises) in the full panel. Classification refers to magnitude, not statistical significance: Large ($|$SDE$| > 0.15$), Moderate ($0.05$--$0.15$), Small ($0.005$--$0.05$), Null ($< 0.005$).
\end{itemize}
}
\end{threeparttable}
\end{table}

